\section{Results}

\subsection{The main result: 200 drives in clear weather}
\begin{plainwords}
The Proposed Bird's-Eye-View Driving Model succeeds in 92.5\% of 200 drives, against 85.0\% and 74.5\% for the two baselines, and 95.0\% for the rule-based Collision-Aware Teacher it learned from. The proposed driver and its teacher cannot be told apart statistically. The image-only driver, which reads the camera picture instead of the grid, succeeds in only \ImOnSuccess{} of the drives.
\end{plainwords}
\explainedfigure{fig_headline_success_clear_weather_200_episodes}{0.97}{Success rate of the driving models on the same 200 test episodes}{Each bar is the share of 200 identical test episodes (scenario seeds 5000 to 5199, clear weather, sparse traffic) in which a model reached the destination without a collision and without leaving the road. The black line is the 95\% Wilson confidence interval, the range of success rates that is statistically compatible with the observed count. Grey bars are rule-based teachers that see the exact simulator state; orange bars are the two baseline models supplied by the teammates; blue bars are models trained in this project.}{The Proposed Bird's-Eye-View Driving Model succeeds in 92.5\% of the episodes. The two baseline models succeed in 85.0\% (trained without traffic) and 74.5\% (trained with traffic). The proposed model is about as good as the Collision-Aware Teacher (95.0\%) whose behaviour it learned; the intervals overlap, so the two cannot be told apart. The Image-Only Driving Model, which reads camera pixels instead of the bird's-eye-view grid, reaches 0.130 (interval 0.090 to 0.184), far below every model that uses the grid.}{Training a model by letting it drive and asking a teacher what it should have done (dataset aggregation) produces a model that matches its teacher and clearly beats few-demonstration behaviour cloning. This figure does not say \emph{which part} of the proposed design (the training method or the map overlay) causes the gain; that needs the additional evaluations listed in the paper-directions section.}{fig:headline}

\explainedfigure{fig_paired_difference_headline_versus_others}{0.97}{Paired difference in success between the proposed model and each other model}{For each comparison the difference (proposed model minus the other model) is computed on the same 200 episodes. The point is the difference; the black line is the 95\% bootstrap interval obtained by resampling episodes 5000 times. A line entirely to the right of the zero line means the proposed model is better with statistical support.}{The proposed model is better than the baseline trained without traffic by 0.075 (interval $+0.035$ to $+0.120$) and better than the baseline trained with traffic by 0.180 (interval $+0.125$ to $+0.240$). Its differences to the Collision-Aware Teacher ($-0.025$), the caution-study control ($+0.030$) and the stock driver model ($+0.040$) have intervals that include or touch zero. The Image-Only Driving Model is far below it.}{The claim that the proposed model beats the two baselines is statistically supported. No claim is made that it beats the teachers or the control. Using paired comparisons (same scenarios for every model) is what makes a difference of 0.075 detectable with 200 episodes.}{fig:paired}


\begin{cautionbox}[What the main result does not show]
It does not separate \emph{why} the proposed driver is better: three things changed together (repeated drive-and-correct teaching instead of five recordings; the map overlay; the choice of three vehicle-state numbers). The plain-grid variant with six numbers was tested on only 10 and 40 drives (0.70 and 0.875 against 0.80 and 0.925), not on 200. Nothing is claimed about dense traffic, where even the teachers fail more than half of the drives. The results are for the camera-free surrogate.
\end{cautionbox}

\subsection{The image-only driver}
The Image-Only Driving Model sees only a front-camera picture (128 by 96 pixels): no grid, no numbers. On the same 200 drives it succeeds in \ImOnSuccess{} (interval \ImOnLow{} to \ImOnHigh{}); \ImOnCollision{} of the drives end in a collision, \ImOnOffroad{} by leaving the road and \ImOnTimeout{} by timeout. Its mean speed is \ImOnSpeed{} metres per second and its estimated energy use \ImOnEnergy{} watt-hours per kilometre.
\begin{implicationbox}
Reading the picture directly is far harder than reading the grid. The grid (a compact, structured summary of the surroundings, produced by the detector and the stereo depth) is what makes the proposed driver possible at this amount of training and data. Whether a camera network trained with far more data could close the gap is not answered here.
\end{implicationbox}

\subsection{Does weather make driving worse?}
\begin{plainwords}
First we measured the damage. In clear weather the object detector finds 72\% of the vehicles in view; in the strongest fog 31\%; in rain of intensity 0.6 only 15\%; in the heaviest rain 1.4\%. The distance estimate for vehicles that \emph{are} found does not get worse, so the damage is almost entirely \emph{missed vehicles}.
Then we tested the drivers with that damage applied. \textbf{Driving success did not fall:} 89--91\% for the proposed driver in all four weather conditions. A pre-registered hypothesis that the damage would hurt a driver trained without it was therefore \emph{not} confirmed.
\end{plainwords}
\explainedfigure{fig_measured_perception_degradation}{0.97}{Measured effect of simulated fog and rain on the camera perception}{Left: the share of vehicles in view that the YOLO object detector finds, per weather condition (vertical lines are 95\% intervals). Right: the share of vehicle pixels at which the stereo-depth algorithm returns a valid distance. Weather conditions are the teammates' image augmentation at the intensity shown (0 = none, 1 = strongest). 40 scenes; 72 vehicle observations in clear weather and 144 in each weather condition.}{In clear weather the detector finds 72\% of the vehicles. In fog of intensity 1.0 it finds 31\%, in rain of intensity 0.6 it finds 15\% and in rain of intensity 1.0 only 1.4\%. Stereo depth stays valid for 88\% of pixels in clear weather and 49\% in the strongest fog, but remains at 67\% to 83\% in rain.}{The perception measurement confirms that the simulated weather is a real degradation of the camera pipeline, mostly through \emph{missed vehicles}. For vehicles that are detected, the distance estimate is not worse. This tells us what a driving model would have to cope with, and explains why the driving results in the next figures are unaffected.}{fig:perception}

\explainedfigure{fig_success_by_weather_calibrated_perception_200_episodes}{0.97}{Success rate in four weather conditions under the measured perception degradation}{Four models were each evaluated on the same 200 episodes in clear weather, fog of intensity 1.0, rain of intensity 0.6 and rain of intensity 1.0. In every weather condition the input of the model was degraded as measured: vehicles can disappear from the bird's-eye-view grid, or be placed at a wrong distance. Black lines are 95\% intervals.}{The proposed model stays between 0.89 and 0.91 in all four conditions; the weather-degradation-trained model between 0.88 and 0.90. The reinforcement-learning-fine-tuned model is lower in every condition (0.765 to 0.815) and the baseline trained with traffic is flat at about 0.75.}{Under the perception degradation that was measured and implemented, weather does not lower driving success in sparse traffic. It follows that a robustness or weather-adaptation claim cannot be made from this experiment, and that a weather-adaptive method would have nothing to adapt to here. The next research step is to find the conditions (denser traffic, harsher loss of vehicles) under which weather begins to matter.}{fig:weather}


\begin{implicationbox}[What this means for the original question]
The question ``can the driver adapt to degraded perception?'' cannot be answered here, because the degradation does not hurt sparse-traffic driving. The research step that follows is to find the conditions in which weather \emph{does} start to matter (denser traffic, more severe loss of vehicles), which is Direction 2 in Section~\ref{sec:g-next}.
\end{implicationbox}

\subsection{Caution-aware driver against the proposed driver, 200 drives in each weather}
On 200 drives per condition the caution-aware driver matches the proposed driver in clear weather (0.890 against 0.925, not significant), fog (0.895 against 0.895) and rain 0.6 (0.890 against 0.900), and is lower in the strongest rain (0.820 against 0.900, $p=0.033$, all extra failures are timeouts). In weather it uses about 35\% less estimated energy (59.7 to 62.7 against 93.0 to 96.5 watt-hours per kilometre). All the extra failures are timeouts: with a 160-second limit instead of 80 seconds (a post-hoc check) the caution-aware driver finishes 0.940, 0.955 and 0.935 in fog 1.0, rain 0.6 and rain 1.0, a tie with the proposed driver (no difference is significant) with fewer collisions. Full table, paired tests and figure: Complete Technical Research Report, section ``Caution-aware model against the proposed model and both baselines''.

\subsection{Three side studies}
\paragraph{Caution when perception is unreliable (40 drives, pilot).} A driver given a ``caution'' number derived from the measured detector reliability slowed from about 8.0 to 5.6--6.2 metres per second in weather and used about 30\% less estimated energy (84.7 to 56--59 watt-hours per kilometre), with more timeouts (7.5--15\%) and no gain in success. This is \emph{learned adaptive behaviour}: the driver was taught by a teacher whose speed and gap depend on caution, and reproduces that dependence by itself, while the control with caution fixed at zero does not slow down. The benefit is not shown (success did not rise; weather did not hurt the proposed driver anyway), and the caution value is supplied from the known weather rather than sensed by the driver.
\paragraph{Which input matters (10 drives, pilot).} The map overlay with the three vehicle-state numbers was nominally the best design (0.80 in sparse and 0.60 in dense traffic); the overlay with no numbers was the worst (0.30 and 0.10). Intervals are very wide.
\paragraph{A wrong position estimate (40 drives, pilot).}
\explainedfigure{fig_success_versus_localisation_error_pose_trained_models}{0.97}{Training with position error, and adding a second lane cue}{The proposed model reads the road from a map drawn into its input, so it needs to know where the car is on the map. Here the position estimate is shifted by a constant lateral distance (horizontal axis, metres; a small heading error grows with it). Success over 40 episodes for the proposed model and four models trained with random position errors; bands are 95\% intervals. Dashed lines: the map is the only source of lane position. Solid green and pink lines: a second, independent and noisy lane measurement (lateral offset and heading) is also given to the model. 40 episodes per point.}{Training with position errors alone does not help (dashed lines, 0.25 to 0.28 at 2 metres of error). Adding a second independent lane cue keeps success between 0.875 and 0.95 for all error sizes up to 2 metres (solid lines).}{Robustness to a wrong position estimate comes from \emph{redundancy} (two sources that can disagree), not from noise augmentation. Caution: the second cue is simulated (simulator truth plus noise), so the result shows that the network can exploit redundancy, not that a real lane detector would be accurate enough.}{fig:posetrained}

